\section*{Overview}

For serious 64-bit primality and factorization tasks, trial division is not enough. The standard contest toolkit is:

\begin{itemize}
  \item Miller-Rabin to decide primality quickly,
  \item Pollard Rho to split a composite number into nontrivial factors.
\end{itemize}

The two techniques are usually used together, not separately.

\section*{Problem-Driven Motivation}

Suppose inputs are arbitrary 64-bit integers and the task asks for:

\begin{itemize}
  \item primality testing for many queries,
  \item full prime factorization,
  \item divisor counts or multiplicative-function values of huge numbers.
\end{itemize}

Trial division up to \(\sqrt{n}\) is too slow once \(n\) reaches around \(10^{18}\). But a full heavyweight number
theory library is unnecessary. What we need is:

\begin{quote}
Fast primality rejection, plus a practical way to split composites.
\end{quote}

\section*{Recognition Pattern}

Reach for this pair when:

\begin{itemize}
  \item numbers fit in 64 bits but are too large for \(\sqrt{n}\) methods,
  \item randomization is acceptable,
  \item the task needs exact factorization, not just small-prime preprocessing.
\end{itemize}

If all numbers are small enough for a sieve or bounded trial division, this machinery is unnecessary.

\section*{Derivation}

\paragraph{Miller-Rabin.}
For odd prime \(n\), write
\[
n-1 = d \cdot 2^s
\]
with \(d\) odd. Then for any base \(a\), the powers
\[
a^d,\ a^{2d},\ a^{4d},\ \dots,\ a^{2^{s-1}d}
\]
must satisfy a very rigid modular pattern. Composites usually violate it, and such a violating base is called a
witness.

So Miller-Rabin is not proving primality directly. It is trying to find a contradiction to primality extremely quickly.

\paragraph{Pollard Rho.}
Once we know \(n\) is composite, we do not need the smallest factor. Any nontrivial factor is enough. Pollard Rho
creates a pseudorandom walk
\[
x_{k+1} = f(x_k) \bmod n
\]
and compares two iterates with Floyd's cycle trick. If the walk collides modulo one hidden prime factor earlier than it
collides modulo the full number, the gcd of the difference with \(n\) leaks a nontrivial factor.

\section*{Worked Problem}

\paragraph{Problem.}
Given many 64-bit integers, return the largest prime factor of each.

\paragraph{Why naive fails.}
Trial division up to \(\sqrt{n}\) is far too slow on worst-case semiprimes near \(10^{18}\).

\paragraph{Step 1: primality test.}
Run deterministic Miller-Rabin for 64-bit integers using a fixed witness set. If \(n\) is prime, stop immediately.

\paragraph{Step 2: split composites.}
If \(n\) is composite, run Pollard Rho until it returns a nontrivial divisor \(d\).

\paragraph{Step 3: recurse.}
Factor \(d\) and \(n/d\) recursively. The recursion eventually bottoms out at primes because every split strictly
reduces the number.

\paragraph{Why this is practical.}
Miller-Rabin prevents wasting Pollard-Rho effort on actual primes. Pollard Rho only needs one factor, not the best
factor, so random splitting is enough.

\section*{Implementation Reasoning}

Three implementation details matter a lot:

\begin{itemize}
  \item use \texttt{\_\_int128} for modular multiplication to avoid overflow,
  \item hard-code the deterministic witness set valid for 64-bit integers,
  \item restart Pollard Rho with a new random constant if one walk fails by returning \(n\).
\end{itemize}

The recursive factorizer is short because Miller-Rabin and Pollard Rho each do one job cleanly.

\section*{Correctness Intuition}

Miller-Rabin is correct because primes satisfy the strong modular squaring pattern, so any witness really proves
compositeness. Pollard Rho is effective because the pseudorandom walk tends to collide modulo hidden prime factors before
it fully cycles modulo \(n\), and that mismatch is exposed by the gcd.

\section*{Complexity Analysis}

Miller-Rabin on 64-bit integers uses a constant number of witnesses, so its cost is essentially a small number of
modular exponentiations. Pollard Rho is randomized, so the cleanest bound is probabilistic, but in practice it is fast
enough for standard contest-sized 64-bit inputs.

\section*{Common Pitfalls}

\begin{itemize}
  \item Overflowing modular multiplication without \texttt{\_\_int128}.
  \item Using an insufficient witness set and silently misclassifying a composite as prime.
  \item Forgetting small cases like \(n < 2\), even numbers, or small primes.
  \item Treating Pollard Rho as deterministic; failed random walks need restarts.
\end{itemize}

\section*{Variants and Failure Modes}

\begin{itemize}
  \item For very small numbers, trial division is simpler.
  \item For larger-than-64-bit integers, the same ideas still work but the implementation is no longer deterministic with
  the same witness set.
  \item Small-prime stripping before Pollard Rho is often a good practical improvement.
\end{itemize}

\section*{Practice Problems}

\begin{itemize}
  \item Factor arbitrary 64-bit integers.
  \item Answer many primality queries on large inputs.
  \item Compute divisor counts or largest prime factors from the returned factorization.
\end{itemize}

\section*{References}

\begin{itemize}
  \item \href{https://cp-algorithms.com/algebra/factorization.html}{cp-algorithms: Integer Factorization}.
  \item \href{https://github.com/kth-competitive-programming/kactl}{KACTL}.
  \item \href{https://codeforces.com/catalog}{Codeforces Catalog}.
\end{itemize}
